% Summary Document Template
% Fill in PLACEHOLDERs and uncomment templates as needed.
% Writing order: Methodology -> Results -> Discussion -> Conclusion -> Appendices -> Introduction (last)

\documentclass[12pt]{article}

% --- Packages ---
\usepackage{lmodern}
\usepackage[T1]{fontenc}
\usepackage{amsmath,amssymb,amsfonts,mathtools}
\usepackage{booktabs}
\usepackage[authoryear]{natbib}
\usepackage{graphicx}
\usepackage{float}
\usepackage{multirow}
\usepackage{subcaption}
\usepackage{hyperref}
\usepackage{xcolor}
\usepackage{enumitem}
\usepackage{setspace}

% --- Page Layout ---
\usepackage[margin=1in]{geometry}
\setlength{\parindent}{0pt}
\setlength{\parskip}{0.8em}

% --- Section Formatting ---
\usepackage{titlesec}
\titleformat{\section}{\normalfont\scshape\centering}{\thesection.}{0.5em}{}
\titleformat{\subsection}{\normalfont\bfseries}{\thesubsection.}{0.5em}{}

% --- Hyperlinks ---
\hypersetup{
  colorlinks   = true,
  urlcolor     = blue,
  linkcolor    = blue,
  citecolor    = black,
}

% --- Captions ---
\usepackage[margin=10pt,font=small,labelfont=bf,labelsep=period]{caption}

% --- Output Directory (relative path from docs/{tier}/{variant_name}/ to output/{tier}/) ---
% Compute: 3 levels up to project root, then into output/{tier}
\newcommand{\outdir}{PLACEHOLDER}

% --- Common Math Shortcuts (add as needed) ---
\newcommand{\E}{\mathbb{E}}
\newcommand{\R}{\mathbb{R}}

\begin{document}

% --- Title Block ---
\begin{center}
{\Large \textbf{PLACEHOLDER: Document Title}}

\medskip

PLACEHOLDER: Authors

\medskip

{\small Internal Report \hspace{1em}|\hspace{1em} \today}
\end{center}


%-----------------------------------------------------%
\section{Introduction}\label{sec:intro}
%-----------------------------------------------------%

% IMPORTANT: Write this section LAST, after all other sections are complete.
% The introduction should:
%   - Overview the approach and its key innovations
%   - Preview the key findings (reference what was found in the results)
%   - Frame findings in terms of the conjectures developed in the Discussion
%   - Set up the narrative arc: make promises the body delivers on

% PLACEHOLDER


%-----------------------------------------------------%
\section{Methodology}\label{sec:method}
%-----------------------------------------------------%

% Brief overview of data sources, key methodological choices, estimation strategy.
% Include key equations from the src/ modules.
% Point to Appendix A for the detailed pipeline walkthrough.

% PLACEHOLDER

% \subsection{Data.}\label{sec:method-data}
% PLACEHOLDER

% \subsection{Estimation.}\label{sec:method-estimation}
% PLACEHOLDER


%-----------------------------------------------------%
\section{Results}\label{sec:results}
%-----------------------------------------------------%

% One subsection per logical group of outputs.
% For each table/figure: include it, write a caption, write 1-3 paragraphs of interpretation.
% Cross-reference related results with \ref{}.

% PLACEHOLDER

% --- Table template ---
% Uncomment and replicate for each table in output/{tier}/tables/{variant_name}/
%
% \begin{table}[H]
%   \centering
%   \caption{PLACEHOLDER: descriptive caption}
%   \label{tab:PLACEHOLDER}
%   \input{\outdir/tables/{variant_name}/PLACEHOLDER}
% \end{table}

% --- Figure template ---
% Uncomment and replicate for each figure in output/{tier}/figures/{variant_name}/
%
% \begin{figure}[H]
%   \centering
%   \includegraphics[width=0.9\textwidth]{\outdir/figures/{variant_name}/PLACEHOLDER.png}
%   \caption{\textbf{PLACEHOLDER: bold title.} PLACEHOLDER: detailed description.}
%   \label{fig:PLACEHOLDER}
% \end{figure}

% --- Subfigure template ---
% Use when multiple related panels should be shown together
%
% \begin{figure}[H]
%   \centering
%   \begin{subfigure}{0.48\textwidth}
%     \includegraphics[width=\textwidth]{\outdir/figures/{variant_name}/PLACEHOLDER.png}
%     \caption{PLACEHOLDER}
%   \end{subfigure}
%   \hfill
%   \begin{subfigure}{0.48\textwidth}
%     \includegraphics[width=\textwidth]{\outdir/figures/{variant_name}/PLACEHOLDER.png}
%     \caption{PLACEHOLDER}
%   \end{subfigure}
%   \caption{PLACEHOLDER: overall caption}
%   \label{fig:PLACEHOLDER}
% \end{figure}

% --- Hand-coded summary table template ---
% Use when synthesizing across multiple outputs (e.g., comparing coefficients across specs)
%
% \begin{table}[H]
%   \centering
%   \caption{PLACEHOLDER: summary title}
%   \label{tab:PLACEHOLDER}
%   \small
%   \begin{tabular}{lcccc}
%     \toprule
%     \textbf{PLACEHOLDER} & \textbf{Col 1} & \textbf{Col 2} & \textbf{Col 3} \\
%     \midrule
%     Row 1 & val & val & val \\
%     Row 2 & val & val & val \\
%     \bottomrule
%   \end{tabular}
%   \medskip
%   {\footnotesize PLACEHOLDER: table notes (significance levels, data sources, etc.)}
% \end{table}


%-----------------------------------------------------%
\section{Discussion}\label{sec:discussion}
%-----------------------------------------------------%

% Synthesize cross-result connections.
% Present conjectures: supporting evidence, complications, what further analysis could resolve each.
% Distinguish well-supported explanations from speculative ones.
% Central finding and implications. Limitations.

% PLACEHOLDER


%-----------------------------------------------------%
\section{Conclusion}\label{sec:conclusion}
%-----------------------------------------------------%

% Main takeaways concisely.
% Open questions and promising directions (from unresolved conjectures).

% PLACEHOLDER


%-----------------------------------------------------%
\appendix
%-----------------------------------------------------%

\section{Detailed Methodology}\label{app:methodology}

% Step-by-step pipeline walkthrough.
% For each stage: what it does, key parameters, which script implements it.
% Include equations where code implements mathematical procedures.

% PLACEHOLDER

% \subsection{Stage 1: PLACEHOLDER.}\label{app:stage1}
% Implemented in \texttt{PLACEHOLDER.py}.
% PLACEHOLDER


\section{Additional Results}\label{app:additional}

% Robustness checks, variant comparisons, additional figures.
% Organize by diagnostic type (e.g., all density plots together).

% PLACEHOLDER


\end{document}
